\section{Introduction}
\label{sec:introduction}

A language model that labels a market regime may be recalling the period it was trained on, or it may be reasoning about the numbers it is shown. Temporal obfuscation separates the two readings: dates and tickers are removed, so that the model has only the values and the stated criteria. The test is only as sound as the numerical series that it shows to the model. This article re-runs the published 30-day regime test on the complete daily series, with a second model on identical inputs, and asks what the published conclusions rest on.

\subsection{Motivation: Data-Definition Drift and Model Retirement}
\label{sec:introduction:motivation}

% src: planning/forensics/2026-10-28_memo_published_inputs.md (raw table reloaded 2026-02-03; options_daily_summary rebuilt 2026-10-08); planning/02_decisions_log.md D-01 (404 model_not_found)
The published run used o4-mini in January 2026. The model no longer answers inference requests (the endpoint returns \ident{model_not_found}). The raw options table was reloaded on 2026-02-03, and the daily GEX table was rebuilt on 2026-10-08 from stored per-contract gamma, with same-day expiries included (Section~\ref{sec:methods:lineage}). A result that depends on a model and on an input definition that can both disappear cannot be re-run, only re-interpreted. The study therefore recovers the exact inputs of the published run from the stored batch file, bridges a second model (gpt-6-luna) to o4-mini on these inputs, and stores every input, prompt and response.

\subsection{Relation to Prior Work}
\label{sec:introduction:relation}

This article is the extended version of the AIAI~2026 conference paper \cite{regan2026regimes}, which introduced 30-day regime detection under temporal obfuscation and built on the single-day study presented at IEEE BigData~2025 \cite{regan2025obfuscation}. Both lines of work are combined in a separate journal article \cite{regan2026jrfm} (Journal of Risk and Financial Management, JRFM).

Relative to all three publications, the new elements are:
\begin{enumerate}[label=(\roman*),leftmargin=*,itemsep=0pt]
\item the provenance and definition of the daily input series;
\item a second model, gpt-6-luna, bridged to o4-mini on identical inputs;
\item the years 2008--2019, the exchange-traded funds (ETFs) SPY (S\&P 500), QQQ (Nasdaq-100) and IWM (Russell 2000), 2026 and an intraday check;
\item overlap-aware inference;
\item an obfuscated-versus-dated rerun of the Phase~3 and~4 baselines; and
\item a disclosed per-symbol scaling for QQQ and IWM.
\end{enumerate}
The recap of the earlier design is limited to Section~\ref{sec:methods}.

\subsection{Research Questions}
\label{sec:introduction:rqs}

% TODO(authors): confirm the wording of RQ1-RQ7 and the mapping of the questions to Sections 4.1-4.11 given after the list.
We address seven questions:
\begin{enumerate}[label=RQ\arabic*.,leftmargin=*]
\item Which daily series were the published inputs, and how do the deterministic labels change on the complete series?
\item Does gpt-6-luna reproduce the o4-mini labels on the exact published inputs?
\item Does the model's regime call follow the input series and the stated criteria?
\item How much of the persistence signal lives in zero-days-to-expiration (0DTE), i.e.\ same-day-expiry, gamma?
\item Do the answers hold for 2008--2019, for QQQ and IWM, and in the post-publication period (2026)?
\item Does a late-session intraday snapshot reproduce the end-of-day label?
\item Does the 2024-versus-2020 contrast of the earlier baseline runs, whose prompts carried calendar dates, persist when the dates are replaced by relative day labels?
\end{enumerate}
RQ1 is answered in Sections~\ref{sec:results:inputs}--\ref{sec:results:replay}, RQ2 in Section~\ref{sec:results:bridge}, RQ3 in Sections~\ref{sec:results:complete} and~\ref{sec:results:fidelity}, RQ4 in Section~\ref{sec:results:ablation}, RQ5 in Sections~\ref{sec:results:deep}--\ref{sec:results:post}, RQ6 in Section~\ref{sec:results:intraday} and RQ7 in Section~\ref{sec:results:dated}.

\subsection{Contributions}
\label{sec:introduction:contributions}

% src: stats.yaml / report.yaml headline (pooled 2024-2025 SPY rule rates: jan2026_phase4a 501 of 501, new_full 60 of 502; replay agreement 1,412 of 1,412);
%      bridge_agreement.yaml bridge[0] (agree_pct 99.6, kappa 0.993); detection_by_year_wide.csv arm=full_spy (luna 613 of 4,671);
%      detection_by_year.csv arm=full_spy scope=all agree_pct 99.67887, kappa 0.985848, n_compared 4671; stats.yaml rates classifier new_ex0dte vs new_full, 2024-2025, rate_pct 14.741036 vs 11.952191;
%      stats.yaml rates classifier new_full_scaled period=all (QQQ 20.958084, IWM 53.919808; 350 of 1,670 and 901 of 1,671); new_full SPY 2024-2025 nonoverlap_30d 3 of 17, block bootstrap 5.009-19.201;
%      report.yaml inputs.detections.rows 36,634; dated vs obfuscated 180 vs 181 of 223 (stats.yaml / decisions log D-22).
This article makes six contributions.
(i)~It reconstructs the published inputs and gives an account of the data they carried. The published 2024 and 2025 values are the call leg of the options chain without the put leg, from 2023-09-29 on. The deterministic labels are recomputed on the complete series: the rule detects 501 of 501 windows (100.0\%) of 2024 and 2025 on the published inputs and 60 of 502 (12.0\%) on the complete series (Section~\ref{sec:results:inputs}, Section~\ref{sec:results:replay}).
(ii)~It analyses the sensitivity of the deterministic labels to the series definition across six daily series (one of them a call-only diagnostic), three assets and 2008--2026. Removing same-day expiries does not restore the published rates, and the fixed thresholds leave QQQ and IWM without a detection until a disclosed scaling is applied (Sections~\ref{sec:results:ablation}--\ref{sec:results:cross}).
(iii)~It bridges o4-mini and gpt-6-luna on identical inputs. The two models give nearly the same verdict on the published windows, and on the complete series gpt-6-luna follows the rule almost everywhere (Sections~\ref{sec:results:bridge}, \ref{sec:results:complete} and~\ref{sec:results:fidelity}).
(iv)~It applies overlap-aware statistics to stride-one windows: the quarter-block bootstrap and a non-overlapping subsample next to the Wilson intervals. Wilson and bootstrap intervals are reported for the pooled rule rate (Section~\ref{sec:methods:statistics} for the method, Section~\ref{sec:results:replay} for the values).
(v)~It stores every input, prompt and response in result tables (36{,}634 answers), from which each table and figure of this article is generated (Section~\ref{sec:methods:storage}).
(vi)~It reruns the Phase~3 and~4 baselines with relative day labels instead of calendar dates. For 2024, gpt-6-luna detects 180 of 223 windows (80.7\%) with relative labels and 181 of 223 (81.2\%) with dates, so the visible dates changed nothing that the data can detect (Section~\ref{sec:results:dated}).

\subsection{Organization}
\label{sec:introduction:organization}

Section~\ref{sec:background} reviews dealer gamma exposure, the validation of LLM reasoning by obfuscation and the construct validity of financial inputs. Section~\ref{sec:methods} defines the daily series, the windows and rule, the model protocol, the experimental arms and the statistics. Section~\ref{sec:results} reports the results, from the lineage of the published inputs to the rerun of the earlier baselines. Section~\ref{sec:discussion} discusses what survives, what the model tracks, the threats to validity and the implications. Section~\ref{sec:conclusion} concludes, and the appendices give the prompt, schema and API configuration (Appendix~\ref{app:prompt}), the series definitions in SQL (Appendix~\ref{app:series}) and the per-year tables (Appendix~\ref{app:tables}).
